% Введение (подключается в annotation.tex и main.tex)
\section{Введение}

Архитектура Transformer~\cite{vaswani2017attention} стала основой
современных больших языковых моделей. Её центральный механизм,
самовнимание (self-attention), требует $O(n^2)$ операций по длине
последовательности $n$. Это породило большое семейство методов
эффективной аппроксимации внимания. Их можно условно разделить на
несколько групп:
\begin{itemize}
  \item \textbf{ядерные и линейные} методы, в которых softmax заменяется
    скалярным произведением признаковых отображений: линейное
    внимание~\cite{katharopoulos2020rnn}, Performer со случайными
    признаками FAVOR+~\cite{choromanski2021performer}, случайные
    фурье-признаки RFA~\cite{peng2021rfa}, SOFT без
    softmax~\cite{lu2021soft};
  \item \textbf{низкоранговые} методы: проекции Linformer по оси
    последовательности~\cite{wang2020linformer}, приближение Нистрёма
    в Nyströmformer~\cite{xiong2021nystrom}, аналитический фреймворк
    A3~\cite{wong2025a3};
  \item \textbf{разреженные} шаблоны внимания: Sparse
    Transformer~\cite{child2019sparse}, хеширование LSH в
    Reformer~\cite{kitaev2020reformer}, комбинация локального,
    глобального и случайного внимания в BigBird~\cite{zaheer2020bigbird};
  \item \textbf{замена внимания} фиксированным или
    синтезированным смешиванием токенов: преобразование Фурье в
    FNet~\cite{leethorp2022fnet}, Synthesizer~\cite{tay2021synthesizer}.
\end{itemize}

Почти все эти методы аппроксимируют механизм внимания \emph{равномерно
по глубине}: одна и та же схема применяется во всех слоях и, как
правило, требует обучения модели с нуля или дообучения. При этом
многие из них по своей природе \emph{стохастичны}: случайные признаки,
случайные хеш-функции, случайные рёбра графа внимания. Из-за этого
ошибка аппроксимации зависит от выборки, а результат вычислений
воспроизводим лишь с точностью до зерна генератора.

Параллельно развивается механистическая интерпретируемость: она
показывает, что слои трансформера неоднородны по своей
функции~\cite{olah2020zoom,rai2024review}. Математический анализ
трансформерных схем~\cite{elhage2021framework} показывает, что модель
без слоёв внимания приближает статистику биграмм. Однослойная модель,
состоящая только из внимания, дополнительно реализует
<<skip-trigram>>-зависимости. Индукционные головы, которые
обеспечивают обучение в контексте, возникают лишь при композиции как
минимум двух слоёв внимания~\cite{olsson2022induction}. Слои прямого
распространения работают как память <<ключ--значение>>, причём ключи
нижних слоёв реагируют на поверхностные шаблоны (конкретные
$n$-граммы), а верхних "--- на семантические~\cite{geva2021ffn}.
Словарное обучение на активациях MLP небольшой модели выделяет
интерпретируемые моносемантические признаки~\cite{bricken2023monosemanticity}.
Известно также, что многоголовое самовнимание способно выражать
свёрточные слои~\cite{cordonnier2020relationship}. Значит, локальные
позиционные паттерны внимания в принципе описываются фиксированными
ядрами. Наконец, теоретический анализ
показывает, что чистое внимание без остаточных связей и MLP
вырождается по рангу дважды экспоненциально с
глубиной~\cite{dong2021attention}. Поэтому остаточный поток играет
ключевую роль при любой послойной замене вычислений.
Отдельно идёт дискуссия о том, являются ли обученные нейросети
<<уже чистыми>> (\emph{already clean}), то есть поддающимися явному
описанию~\cite{cammarata_clean,gao_notclean,lissack_clean,ibm_clean}.
Cammarata~\cite{cammarata_clean} утверждает, что видимая
полисемантичность нейронов в основном объясняется не суперпозицией, а
<<рекурсивным обобщением>> связанных понятий, и по предварительной
оценке около 90\% <<полюсов>> нейронов Qwen3-32B (положительная или отрицательная сторона активации) имеют одну основную функцию.
Gao~\cite{gao_notclean} возражает, что совпадения активирующих примеров с
шаблоном недостаточно, и предлагает строгий критерий: объяснение
верно, если подстановка смоделированных по нему активаций не
увеличивает KL-дивергенцию выходного распределения модели. Именно этот
критерий "--- KL при подстановке "--- мы используем как основную
метрику замены слоёв.

Эти наблюдения подсказывают иную постановку задачи. Если первые слои
LLM реализуют в основном локальные, <<поверхностные>> вычисления, то
их можно заменить \emph{явными детерминированными алгоритмами} вместо
единообразной аппроксимации внимания во всей сети. Под
детерминированным алгоритмом мы понимаем процедуру, которая
\begin{enumerate}
  \item не использует случайную выборку на этапе вывода;
  \item строится по весам обученной модели и/или статистике корпуса
    в замкнутом виде или явной процедурой, без градиентного
    дообучения;
  \item даёт одинаковый выход на одинаковом входе и допускает
    аналитическую оценку ошибки.
\end{enumerate}
Примеры таких алгоритмов: фиксированные разреженные шаблоны внимания,
позиционные ядра вместо голов внимания, таблицы поиска по токенам,
аналитические низкоранговые разложения матриц $W_{QK}$, $W_{OV}$ и
весов MLP~\cite{wong2025a3,yuan2023asvd,wang2025svdllm}, линейные
карты, подогнанные гребневой регрессией.

\paragraph{Постановка задачи.}
Пусть $f = f_L \circ \dots \circ f_1 \circ E$ "--- языковая модель из
$L$ блоков с матрицей эмбеддингов $E$, а $h_k(x)$ "--- состояние
остаточного потока после $k$-го блока на входе $x$. Требуется
построить детерминированный алгоритм $A_k$, такой что
\[
  \tilde f = f_L \circ \dots \circ f_{k+1} \circ A_k,
  \qquad
  \mathbb{E}_{x}\,
  \mathrm{KL}\bigl(p_f(\cdot \mid x)\,\|\,p_{\tilde f}(\cdot \mid x)\bigr)
  \le \varepsilon,
\]
и определить максимальную глубину $k^{*}(\varepsilon)$, для которой
это возможно при ограничениях на вычислительную сложность и память
алгоритма $A_k$.

\paragraph{Вклад работы.}
\begin{enumerate}
  \item Формализация задачи детерминированной аппроксимации первых
    $k$ слоёв LLM с критерием качества <<KL при подстановке>> и
    протокол, в котором исходная и изменённая модели сравниваются
    в одном процессе в fp32 (шум bf16 сам по себе даёт KL$\approx10^{-3}$).
  \item Адаптация к каузальному декодеру с GQA, QK-нормализацией и RoPE
    шести методов эффективного внимания и пяти методов сжатия MLP из
    литературы; разбор того, почему Linformer, Nyströmformer, FNet,
    SOFT и Reformer неприменимы без переобучения.
  \item Новые детерминированные методы: факторизация внимания
    <<позиционное ядро $\times$ значимость ключевого токена>> (O4),
    типизированная замена голов (O1), ридж-свёртка вместо
    внимания (O2), аддитивная $n$-граммная замена блоков (O3),
    линеаризация MLP с точными нейронами (M-O1), табличный (M-O2) и
    кусочно-линейный (M-O3) MLP. Все наши методы калибруются
    последовательно с компенсацией накопленного дрейфа.
  \item Послойный профиль Qwen3-0.6B и 1.7B, выделяющий режим <<до стока
    внимания>> и задающий глубину $k_{\mathrm{sink}}$, и анализ
    неприменимости случайных признаков к первым слоям.
  \item Сравнение по KL, perplexity и семи бенчмаркам на двух
    размерах модели, кривые <<качество -- бюджет вычислений>>,
    доверительные интервалы и дистилляционный контроль. Итог: внимание
    первых трёх слоёв заменяется окном со стоком за 1.5--1.8\,\% PPL без
    потерь на бенчмарках; MLP первых слоёв "--- лишь частично, лучший
    метод "--- наш предиктор активных нейронов M-O4 (4.7\,\% на 0.6B,
    1.7\,\% на 1.7B при четверти вычислений MLP).
  \item Ловушка калибровки в замкнутой форме: единичные токены с
    массивной активацией доминируют в задаче наименьших квадратов и
    обрушивают модель на слое, где возникает сток внимания. Мы даём
    формулу влияния строки, воспроизводим эффект дозой отравленных окон,
    устраняем его исключением строк-выбросов и показываем, какие
    опубликованные методы ему не подвержены и почему.
  \item Нейроны стока: массивную активацию создают один--пять нейронов
    MLP, без которых модель не работает; их сохранение "--- обязательное
    условие любого отбора нейронов в ранних слоях.
  \item Карта заменяемости по глубине и априорный критерий: дешёвые
    статистики слоя предсказывают KL его замены с ранговой корреляцией
    0.9--0.97, воспроизводятся на второй модели и позволяют планировать
    замены на непросмотренной модели в 10--15 раз точнее замены префикса.
\end{enumerate}

\paragraph{Структура статьи.}
Раздел~2 содержит обзор литературы. В разделе~3 описаны постановка,
адаптированные методы из литературы и предложенные методы. Раздел~4
описывает экспериментальную установку, раздел~5 "--- результаты, раздел~6
"--- обсуждение и ограничения.
